\subsection{Cartier indices and specialization through a singular quotient}
\label{subsec:noncartier-specialization}

This subsection applies the cyclic quotient of
\citet[Example~5.10, pp.~24--25]{LiedtkeMatsumoto2018} to specialization through a
non-Cartier vertical divisor. Their construction has a regular covering
surface and a normal singular quotient with nonreduced special fibre.
We use the same construction over a complete discrete valuation ring
containing $\mu_p$, and compute its Cartier index and an explicit
perfect-residue arc. The purpose is to identify the extra local
information needed to pass from a divisorial ramification bound to a
bound on every specialization.
For broader context on wild surface quotients and their divisor and
discriminant groups, see \citet{Lorenzini2013WildQuotients} and
\citet{LorenziniSchroer2023WildQuotients}.

\begin{lemma}[The local Cartier-index constraint]
\label{lem:cartier-index-contact}
Let $R$ be a normal Noetherian local domain of dimension two and mixed
characteristic $(0,p)$. Suppose that
\[
                   \operatorname{div}_R(p)=N D
\]
for a single height-one prime divisor $D$ and an integer $N>0$.
Let $i$ be the least positive integer for which $iD$ is Cartier at the
closed point. Then $i$ exists and divides $N$. If
$\phi:R\longrightarrow\mathcal O_K$ is a local map to a complete
mixed-characteristic discrete valuation ring, with $p\ne0$ and
$e_K=v_K(p)$, then
\begin{equation}
               \frac{N}{\gcd(N,e_K)}\mid i.
\label{eq:cartier-specialization-divisibility}
\end{equation}
In particular, $N=1300$ and $e_K=260$ require $5\mid i$.
More generally, if $N=52\cdot5^j$ with $j\ge2$, an index-$260$
specialization requires $5^{j-1}\mid i$.
\end{lemma}

\begin{proof}
The integers $a$ for which $aD$ is locally principal form a subgroup of
$\mathbf Z$ containing $N$. Since every invertible module over a local
ring is free, its positive generator is the Cartier index $i$, and
$i\mid N$.
Choose $f\in\operatorname{Frac}(R)$ with
$\operatorname{div}_R(f)=iD$. Effectivity and the height-one intersection
property for a normal Noetherian domain give $f\in R$; see
\cite[Lemma~10.157.6]{StacksHeightOneIntersection}.
The rational function $p/f^{N/i}$ and its inverse have nonnegative
valuation at every height-one prime. The same intersection property
therefore gives
\[
                      p=u f^{N/i},\qquad u\in R^\times.
\]
The image of $u$ under $\phi$ is a unit, and the image of $f$ is nonzero.
Taking valuations yields
\[
                      e_K=(N/i)v_K(\phi(f)).
\]
Thus $N\mid i e_K$, which is equivalent to
\eqref{eq:cartier-specialization-divisibility}.
The two numerical consequences follow by taking the indicated greatest
common divisors.
\end{proof}

If the reduced vertical divisor is Cartier, the same proof gives
$e_K\ge N$ for a local specialization through it. In general it gives
only $e_K\ge N/i$. The following example attains this latter bound.

\begin{example}[The Liedtke--Matsumoto quotient at index $260$]
\label{ex:noncartier-quotient-specialization}
There is a complete field $K$ of mixed characteristic $(0,5)$ with
algebraically closed residue field $k$, $\mu_5\subset K$, and $e_K=260$,
and a finite quotient $A=B^{C_5}$ of a regular local surface $B$, such
that $A$ is a normal local surface and the following hold.
The generic-fibre cover is finite \'etale of degree five.
The unique reduced vertical prime $D$ has residue field $k(Y)$,
normalized value $v_D(5)=1300$, and Cartier index five.
The cover of divisorial valuation rings has ramification index one and
purely inseparable residue degree five.
Nevertheless there is a local map $A\to\mathcal O_K$, not contained in
$D$, for which
\[
            \operatorname{length}_{\mathcal O_K}
            (\mathcal O_K/I_D\mathcal O_K)=1.
\]
Here $I_D$ denotes the prime ideal of the reduced divisor.
The generic pullback of the cover is a connected degree-five field
extension of $K$. Its normalization has absolute index $1300$ and
residue field $k$.
\end{example}

\begin{proof}
We first give the same cyclic construction for any prime $p$ and any complete
mixed-characteristic discrete valuation ring $\mathcal O$ with
uniformizer $\pi$, residue field $k$ of characteristic $p$, and a
primitive $p$th root of unity $\zeta\in\mathcal O$.
Put $e_0=v_{\mathcal O}(p)$ and
\[
 \mathcal O'=\mathcal O[t]/(t^p-\pi),\qquad
 B_0=\mathcal O'[U],\qquad B=(B_0)_{(t,U)}.
\]
Eisenstein's criterion shows that $\mathcal O'$ is a complete discrete
valuation ring with uniformizer $t$ and residue field $k$.
Thus $B$ is a regular local ring of dimension two. The formulas
\[
                 \sigma(t)=\zeta t,\qquad
                 \sigma(U)=\zeta^{-1}U
\]
define a faithful action of the constant cyclic group $C_p$, fixing
$\mathcal O$. Set $X=tU$ and $Y=U^p$. We claim that
\begin{equation}
 A_0:=B_0^{C_p}=\mathcal O[X,Y]/(X^p-\pi Y),\qquad
 A:=B^{C_p}=(A_0)_{(\pi,X,Y)}.
\label{eq:noncartier-cyclic-ring}
\end{equation}

Indeed, every element of $B_0$ has a unique finite expansion in monomials
$t^aU^b$, with $0\le a<p$, $b\ge0$, and coefficients in $\mathcal O$.
The monomial has character $\zeta^{a-b}$. Since the coefficient domain
has characteristic zero, invariance forces $b=a+pj$ for an integer
$j\ge0$, giving $t^aU^b=X^aY^j$.
Conversely these monomials are invariant. They show that $A_0$ is free
on $1,X,\ldots,X^{p-1}$ over $\mathcal O[Y]$, with the single relation
$X^p=\pi Y$. The ring $B_0$ is finite over $A_0$, generated as a module
by $t^aU^b$ with $0\le a,b<p$.

There is only one prime of $B_0$ above $(\pi,X,Y)$, because its closed
fibre has ring
\[
                         k[t,U]/(t^p,tU,U^p).
\]
Localization over $A_0$ therefore gives $B$.
To verify invariants after localization directly, write a fixed
element as $f/s$ with $s\notin(t,U)$. The norm
$\prod_{g\in C_p}g(s)$ is invariant and is outside $(\pi,X,Y)$.
Multiplication by this norm gives an invariant numerator in $B_0$,
hence one in $A_0$. This proves~\eqref{eq:noncartier-cyclic-ring}.

The ring $A$ is a Noetherian local domain of dimension two. It is normal:
an element of $\operatorname{Frac}(A)$ integral over $A$ is integral over
$B$, lies in the normal ring $B$, and is fixed, so belongs to $A$.
Here normality of regular rings is the standard result
\cite[Lemma~10.157.5]{StacksRegularNormal}.
After inverting $p$, both $\pi$ and $t$ are units and $U=X/t$, so
\[
                    B[1/p]=A[1/p][t]/(t^p-\pi).
\]
The degree is $p$: the $p$ distinct elements $\zeta^a t$ are conjugate,
and $t^p=\pi$. The derivative $pt^{p-1}$ is a unit, proving finite
\'etaleness of the generic-fibre cover.

Modulo $\pi$, the equation of $A$ is $X^p=0$. Hence its unique reduced
vertical prime is $D=(\pi,X)$, with residue field $k(Y)$.
At $D$, the element $Y$ is a unit and $\pi=X^p/Y$.
The discrete valuation ring $A_D$ has uniformizer $X$, and
\begin{equation}
       v_D(\pi)=p,\qquad v_D(p)=p e_0,
       \qquad\operatorname{div}_A(\pi)=pD.
\label{eq:noncartier-vertical-values}
\end{equation}
The corresponding prime of $B$ is $E=(t)$, where $U$ is a unit.
Thus $v_E(X)=v_E(tU)=1=v_D(X)$, giving relative ramification index one.
The residue extension is $k(Y)\subset k(U)$ with $Y=U^p$, purely
inseparable of degree $p$.

The Cartier index of $D$ divides $p$ by
\eqref{eq:noncartier-vertical-values}. It is not one.
The classes of $\pi,X,Y$ are linearly independent in the cotangent
space of $A$, because $X^p-\pi Y$ belongs to the square of the maximal
ideal of the ambient regular local ring
$\mathcal O[X,Y]_{(\pi,X,Y)}$.
Consequently $I_D=(\pi,X)$ has two independent classes modulo
$\mathfrak m_A I_D$: this latter ideal is contained in $\mathfrak m_A^2$.
The ideal $I_D$ is therefore not principal. In a normal domain it is
the divisorial ideal $\mathcal O(-D)$, so a Cartier divisor $D$ would
make it invertible, hence principal over the local ring. This is a
contradiction. Since $p$ is prime, the Cartier index is exactly $p$.

Now define
\begin{equation}
                  A\longrightarrow\mathcal O,
                  \qquad X\longmapsto\pi,\quad
                         Y\longmapsto\pi^{p-1}.
\label{eq:noncartier-specialization-map}
\end{equation}
The relation is satisfied, and every localized denominator maps to a
unit. The map is local, is not contained in the vertical divisor,
and takes $I_D$ to $(\pi)$, proving the stated intersection length one.
Its target has absolute index $e_0$, while the vertical valuation has
index $pe_0$.
On generic fibres the pullback is
$\operatorname{Frac}(\mathcal O)[t]/(t^p-\pi)$, with
$U=\pi/t=t^{p-1}$. It is an Eisenstein degree-$p$ field extension.
The normalization of its dominating integral component is
$\mathcal O'$, of absolute index $pe_0$ and residue field $k$.
This assertion does not require the whole closed-point tensor product
to be flat over $\mathcal O$.

For the numerical example take $p=5$ and $k=\overline{\mathbf F}_5$.
Start with $\operatorname{Frac}W(k)$, adjoin a primitive fifth root
$\zeta$, and then adjoin $\pi$ with $\pi^{65}=\zeta-1$.
The polynomial
\[
                   \Phi_5(1+Z)=Z^4+5Z^3+10Z^2+10Z+5
\]
is Eisenstein at $5$, so the first extension has degree four and
uniformizer $\zeta-1$. The second polynomial is Eisenstein of degree
$65$. The resulting complete field $K$ has residue $k$, contains
$\mu_5$, and has $e_K=4\cdot65=260$.
Applying the construction with $\mathcal O=\mathcal O_K$ gives
$v_D(5)=5\cdot260=1300$ and every asserted property.
\end{proof}

The example shows that normality, a finite constant-group quotient,
and generic \'etale structure do not by themselves
promote the specified divisorial bound to all perfect-residue curves.
Its reduced vertical prime has precisely the Cartier-index denominator
allowed by Lemma~\ref{lem:cartier-index-contact}.
This ring and curve are auxiliary examples: they have not been identified
with the full Drinfeld-level or nonsplit Cartan quotient, and they supply
no pair of height-two groups with unequal special characters and
isomorphic $5^4$-torsion.
For that actual comparison, it remains to decide whether its finite-level
marking factors through the exceptional-torus chart of
Proposition~\ref{prop:cartan-torus-regularity}, or through another quotient.
That proposition resolves the local divisor question on its stated charts;
it does not identify all curves compatible with the torsion isomorphism.
